\documentclass[12pt,a4paper]{article}

% ---- Encoding & Font ----
\usepackage{iftex}
\ifPDFTeX
  \usepackage[utf8]{inputenc}
  \usepackage[T1]{fontenc}
\else
  \usepackage{fontspec}  % XeTeX/LuaTeX (tectonic): Unicode nativo (·, —, ¿, ª, §)
\fi
\usepackage[spanish]{babel}
\usepackage{csquotes}

% ---- Page layout ----
\usepackage[top=2.5cm, bottom=2.5cm, left=2.5cm, right=2.5cm]{geometry}
\usepackage{setspace}
\onehalfspacing

% ---- Math ----
\usepackage{amsmath, amssymb, amsthm}
\usepackage{mathtools}
\usepackage{physics}
\usepackage{esint}  % \oiint

% ---- Graphics ----
\usepackage{graphicx}
\usepackage{tikz}
\usepackage{float}
\usepackage{caption}

% ---- Tables ----
\usepackage{array}
\usepackage{booktabs}
\usepackage{tabularx}
\usepackage{colortbl}

% ---- Colours ----
\usepackage{xcolor}
\definecolor{notebg}{HTML}{FFF3CD}
\definecolor{noteborder}{HTML}{FFC107}
\definecolor{boxred}{HTML}{F8D7DA}
\definecolor{boxredborder}{HTML}{F5C6CB}

% ---- Boxes ----
\usepackage{mdframed}
\usepackage{tcolorbox}
\tcbuselibrary{most}

\newtcolorbox{keybox}[1][]{
  colback=blue!5,
  colframe=blue!70!black,
  arc=4pt,
  boxrule=1pt,
  fonttitle=\bfseries,
  title={#1}
}

\newtcolorbox{warnbox}[1][]{
  colback=boxred,
  colframe=boxredborder,
  coltitle=black!75,
  arc=4pt,
  boxrule=1pt,
  fonttitle=\bfseries,
  title={#1}
}

\newtcolorbox{notebox}[1][]{
  colback=notebg,
  colframe=noteborder,
  coltitle=black!75,
  arc=4pt,
  boxrule=1pt,
  fonttitle=\bfseries,
  title={#1}
}

% ---- Diagramas (gráficas y esquemas) ----
\usepackage{pgfplots}
\pgfplotsset{compat=1.18}
\usetikzlibrary{babel}  % babel español activa `>`; esto evita romper las flechas `->`
% courses/ElecMag2024/Clases/assets/tikzstyles.tex
% ---------------------------------------------------------------------------
% Estilos compartidos para los diagramas del curso Electromagnetismo (2024).
% Fuente única del "look" visual (colores, grosores, escalas) para que todas
% las figuras (TikZ / pgfplots / CircuiTikz) sean consistentes entre clases.
%
% Es una COPIA de courses/Fisica3-2015/Clases/assets/tikzstyles.tex: mismo
% dominio, misma paleta. Copia y no symlink para que cada curso sea
% autocontenido y la paleta de Física III pueda evolucionar aparte.
%
% Uso: la clase debe cargar antes `tikz` y `pgfplots` (y `circuitikz` si la
%      clase tiene circuitos), y luego % courses/ElecMag2024/Clases/assets/tikzstyles.tex
% ---------------------------------------------------------------------------
% Estilos compartidos para los diagramas del curso Electromagnetismo (2024).
% Fuente única del "look" visual (colores, grosores, escalas) para que todas
% las figuras (TikZ / pgfplots / CircuiTikz) sean consistentes entre clases.
%
% Es una COPIA de courses/Fisica3-2015/Clases/assets/tikzstyles.tex: mismo
% dominio, misma paleta. Copia y no symlink para que cada curso sea
% autocontenido y la paleta de Física III pueda evolucionar aparte.
%
% Uso: la clase debe cargar antes `tikz` y `pgfplots` (y `circuitikz` si la
%      clase tiene circuitos), y luego % courses/ElecMag2024/Clases/assets/tikzstyles.tex
% ---------------------------------------------------------------------------
% Estilos compartidos para los diagramas del curso Electromagnetismo (2024).
% Fuente única del "look" visual (colores, grosores, escalas) para que todas
% las figuras (TikZ / pgfplots / CircuiTikz) sean consistentes entre clases.
%
% Es una COPIA de courses/Fisica3-2015/Clases/assets/tikzstyles.tex: mismo
% dominio, misma paleta. Copia y no symlink para que cada curso sea
% autocontenido y la paleta de Física III pueda evolucionar aparte.
%
% Uso: la clase debe cargar antes `tikz` y `pgfplots` (y `circuitikz` si la
%      clase tiene circuitos), y luego \input{../assets/tikzstyles.tex}
% (compilar desde el directorio de la clase para que la ruta relativa resuelva).
% ---------------------------------------------------------------------------

% ---- Paleta (alineada con las cajas del preámbulo: keybox/notebox/warnbox) ----
\definecolor{figblue}{HTML}{1F4E79}   % azul principal (líneas, keybox)
\definecolor{figred}{HTML}{C0392B}    % rojo de énfasis (warnbox)
\definecolor{figamber}{HTML}{B8860B}  % ámbar (notebox)
\definecolor{figgray}{HTML}{555555}   % auxiliares, ejes, guías

% ---- CircuiTikz: escala y grosor de los componentes ----
% Guarda \ifdefined: la electrostática (Clases 1-14) no carga `circuitikz`, y
% sin la guarda el \input revienta con `Undefined control sequence \ctikzset`.
% Cuando el curso llegue a circuitos, el bloque se activa solo.
\ifdefined\ctikzset
  \ctikzset{
    bipoles/thickness=1,
    resistors/scale=0.9,
    capacitors/scale=0.9,
  }
\fi

% ---- TikZ: estilos reutilizables ----
% Nota: las puntas se declaran como `-{Latex[...]}` y no como `->`. La forma
% con llaves NO contiene un `>` literal, así que es inmune al `>` activo que
% introduce babel español (ver CLAUDE.md §7.3.3).
\usetikzlibrary{arrows.meta,calc,angles,quotes}

\tikzset{
  figline/.style={figblue, line width=0.9pt},
  figaux/.style={figgray, dashed, line width=0.6pt},
  % vectores
  figvec/.style={figblue, line width=0.9pt, -{Latex[length=2.2mm,width=1.6mm]}},
  figvecres/.style={figred, line width=1.3pt, -{Latex[length=2.6mm,width=1.9mm]}},
  figvecaux/.style={figgray, line width=0.7pt, -{Latex[length=1.8mm,width=1.3mm]}},
  figaxis/.style={figgray, line width=0.6pt, -{Latex[length=2mm,width=1.4mm]}},
  % cargas puntuales
  figcarga/.style={circle, inner sep=0pt, minimum size=5.4pt, draw=none},
  figpos/.style={figcarga, fill=figred},
  figneg/.style={figcarga, fill=figblue},
  figneutra/.style={figcarga, fill=figgray},
  % etiquetas
  figlbl/.style={font=\small},
  figlblaux/.style={font=\footnotesize, figgray},
}

% ---- pgfplots: estilo base de las gráficas ----
\pgfplotsset{
  emfig/.style={
    width=11cm, height=6.5cm,
    axis lines=left,
    xlabel style={font=\small},
    ylabel style={font=\small},
    tick label style={font=\footnotesize},
    every axis plot/.append style={line width=1pt},
    grid=major,
    grid style={figgray!20},
    legend cell align=left,
    legend style={font=\footnotesize, draw=figgray!40},
  },
}

% (compilar desde el directorio de la clase para que la ruta relativa resuelva).
% ---------------------------------------------------------------------------

% ---- Paleta (alineada con las cajas del preámbulo: keybox/notebox/warnbox) ----
\definecolor{figblue}{HTML}{1F4E79}   % azul principal (líneas, keybox)
\definecolor{figred}{HTML}{C0392B}    % rojo de énfasis (warnbox)
\definecolor{figamber}{HTML}{B8860B}  % ámbar (notebox)
\definecolor{figgray}{HTML}{555555}   % auxiliares, ejes, guías

% ---- CircuiTikz: escala y grosor de los componentes ----
% Guarda \ifdefined: la electrostática (Clases 1-14) no carga `circuitikz`, y
% sin la guarda el \input revienta con `Undefined control sequence \ctikzset`.
% Cuando el curso llegue a circuitos, el bloque se activa solo.
\ifdefined\ctikzset
  \ctikzset{
    bipoles/thickness=1,
    resistors/scale=0.9,
    capacitors/scale=0.9,
  }
\fi

% ---- TikZ: estilos reutilizables ----
% Nota: las puntas se declaran como `-{Latex[...]}` y no como `->`. La forma
% con llaves NO contiene un `>` literal, así que es inmune al `>` activo que
% introduce babel español (ver CLAUDE.md §7.3.3).
\usetikzlibrary{arrows.meta,calc,angles,quotes}

\tikzset{
  figline/.style={figblue, line width=0.9pt},
  figaux/.style={figgray, dashed, line width=0.6pt},
  % vectores
  figvec/.style={figblue, line width=0.9pt, -{Latex[length=2.2mm,width=1.6mm]}},
  figvecres/.style={figred, line width=1.3pt, -{Latex[length=2.6mm,width=1.9mm]}},
  figvecaux/.style={figgray, line width=0.7pt, -{Latex[length=1.8mm,width=1.3mm]}},
  figaxis/.style={figgray, line width=0.6pt, -{Latex[length=2mm,width=1.4mm]}},
  % cargas puntuales
  figcarga/.style={circle, inner sep=0pt, minimum size=5.4pt, draw=none},
  figpos/.style={figcarga, fill=figred},
  figneg/.style={figcarga, fill=figblue},
  figneutra/.style={figcarga, fill=figgray},
  % etiquetas
  figlbl/.style={font=\small},
  figlblaux/.style={font=\footnotesize, figgray},
}

% ---- pgfplots: estilo base de las gráficas ----
\pgfplotsset{
  emfig/.style={
    width=11cm, height=6.5cm,
    axis lines=left,
    xlabel style={font=\small},
    ylabel style={font=\small},
    tick label style={font=\footnotesize},
    every axis plot/.append style={line width=1pt},
    grid=major,
    grid style={figgray!20},
    legend cell align=left,
    legend style={font=\footnotesize, draw=figgray!40},
  },
}

% (compilar desde el directorio de la clase para que la ruta relativa resuelva).
% ---------------------------------------------------------------------------

% ---- Paleta (alineada con las cajas del preámbulo: keybox/notebox/warnbox) ----
\definecolor{figblue}{HTML}{1F4E79}   % azul principal (líneas, keybox)
\definecolor{figred}{HTML}{C0392B}    % rojo de énfasis (warnbox)
\definecolor{figamber}{HTML}{B8860B}  % ámbar (notebox)
\definecolor{figgray}{HTML}{555555}   % auxiliares, ejes, guías

% ---- CircuiTikz: escala y grosor de los componentes ----
% Guarda \ifdefined: la electrostática (Clases 1-14) no carga `circuitikz`, y
% sin la guarda el \input revienta con `Undefined control sequence \ctikzset`.
% Cuando el curso llegue a circuitos, el bloque se activa solo.
\ifdefined\ctikzset
  \ctikzset{
    bipoles/thickness=1,
    resistors/scale=0.9,
    capacitors/scale=0.9,
  }
\fi

% ---- TikZ: estilos reutilizables ----
% Nota: las puntas se declaran como `-{Latex[...]}` y no como `->`. La forma
% con llaves NO contiene un `>` literal, así que es inmune al `>` activo que
% introduce babel español (ver CLAUDE.md §7.3.3).
\usetikzlibrary{arrows.meta,calc,angles,quotes}

\tikzset{
  figline/.style={figblue, line width=0.9pt},
  figaux/.style={figgray, dashed, line width=0.6pt},
  % vectores
  figvec/.style={figblue, line width=0.9pt, -{Latex[length=2.2mm,width=1.6mm]}},
  figvecres/.style={figred, line width=1.3pt, -{Latex[length=2.6mm,width=1.9mm]}},
  figvecaux/.style={figgray, line width=0.7pt, -{Latex[length=1.8mm,width=1.3mm]}},
  figaxis/.style={figgray, line width=0.6pt, -{Latex[length=2mm,width=1.4mm]}},
  % cargas puntuales
  figcarga/.style={circle, inner sep=0pt, minimum size=5.4pt, draw=none},
  figpos/.style={figcarga, fill=figred},
  figneg/.style={figcarga, fill=figblue},
  figneutra/.style={figcarga, fill=figgray},
  % etiquetas
  figlbl/.style={font=\small},
  figlblaux/.style={font=\footnotesize, figgray},
}

% ---- pgfplots: estilo base de las gráficas ----
\pgfplotsset{
  emfig/.style={
    width=11cm, height=6.5cm,
    axis lines=left,
    xlabel style={font=\small},
    ylabel style={font=\small},
    tick label style={font=\footnotesize},
    every axis plot/.append style={line width=1pt},
    grid=major,
    grid style={figgray!20},
    legend cell align=left,
    legend style={font=\footnotesize, draw=figgray!40},
  },
}


% ---- Hyperlinks & TOC ----
\usepackage[hidelinks]{hyperref}
\usepackage{bookmark}

% ---- Enumerate ----
\usepackage{enumitem}

% ---- Miscellaneous ----
\usepackage{icomma}
\usepackage{siunitx}
\usepackage{textcomp}
\usepackage[bottom]{footmisc}

% ---- Title ----
\title{\textbf{Clase 7: Polinomios de Legendre y separación de variables en esféricas} \\[0.3em]
        \large{La esfera conductora en un campo uniforme}}
\author{Transcripto y expandido --- Curso Electromagnetismo (OpenFING)}
\date{}

\begin{document}

\maketitle
\thispagestyle{empty}
\newpage

\tableofcontents
\newpage


\section{Dónde estamos}

El conjunto de soluciones de una ecuación diferencial lineal \textbf{homogénea}
---incluidas las que son en derivadas parciales, como Laplace--- forma un espacio
vectorial. Sobre esa base, la clase anterior aplicó el \textbf{método de
separación de variables} a la ecuación de Laplace en coordenadas esféricas con
simetría azimutal, buscando soluciones de la forma $\phi(r,\theta) = Z(r)\,P(\theta)$
---no la solución general, sino \textbf{una base} de soluciones---. Eso separó el
problema en dos ecuaciones diferenciales \textbf{ordinarias}:

\[
\dv{r}\big(r^2 Z'(r)\big) = K\,Z(r)
\qquad\qquad
\dv{\theta}\big(\sin\theta\,P'(\theta)\big) + K\sin\theta\,P(\theta) = 0
\]

y el cambio $u = \cos\theta$ llevó la angular a la \textbf{ecuación de Legendre}:

\[
(1-u^2)\,P''(u) - 2u\,P'(u) + K\,P(u) = 0
\]

con el interés puesto en $\theta\in[0,\pi]$, es decir $u\in[-1,1]$.

\begin{warnbox}[El obstáculo]
Es lineal y homogénea, pero \textbf{no de coeficientes constantes}. Con
coeficientes constantes las soluciones son exponenciales ---o combinaciones
complejas de ellas, que dan senos y cosenos---, y se hallan con el polinomio
característico: es la ecuación del resorte, la carga y descarga de un
condensador. Acá los coeficientes dependen de $u$, así que ese camino no existe.
\end{warnbox}

\section{Resolver Legendre por serie de potencias}

\subsection{El método}

El truco es \textbf{buscar la solución como serie de potencias} alrededor de un
punto cómodo. El intervalo de interés es $[-1,1]$, cuyo punto natural del medio
es $u=0$. En lugar de calcular la solución, se calcula su \textbf{desarrollo de
Taylor} alrededor de $u=0$.

\begin{notebox}[Es dos cosas a la vez]
Como \textbf{método numérico} es muy práctico: se calculan términos sucesivos
hasta donde llegue la paciencia y se reemplaza la función desconocida por su
polinomio. Pero acá se va más lejos: se construye la \textbf{serie de Taylor}
completa y se estudia si converge. Ahí aparece el arte del asunto ---la serie
puede converger o no, portarse bien o mal---.
\end{notebox}

Se propone entonces $P(u) = \sum_{m=0}^{\infty} a_m\,u^{m}$ y se deriva término a
término:

\[
P'(u) = \sum_{m=0}^{\infty} a_m\,m\,u^{m-1}
\qquad
P''(u) = \sum_{m=0}^{\infty} a_m\,m(m-1)\,u^{m-2}
\]

\textbf{Reindexado.} En $P''$ los términos $m=0$ y $m=1$ son nulos, así que la
suma arranca realmente en $m=2$. Definiendo $m' = m-2$ y volviendo a llamar $m$
al índice ---es un \textbf{índice mudo}, como en un cambio de variable de una
integral---:

\[
P''(u) = \sum_{m=0}^{\infty} a_{m+2}\,(m+2)(m+1)\,u^{m}
\]

que ya tiene la misma potencia $u^m$ que los demás términos.

\subsection{La relación de recurrencia}

Sustituyendo los cuatro términos de la ecuación, todos escritos con $u^m$:

\begin{table}[H]
  \centering
  \begin{tabular}{@{}ll@{}}
    \toprule
    \textbf{Término} & \textbf{Queda} \\
    \midrule
    $P''$      & $\sum_m a_{m+2}(m+2)(m+1)\,u^m$ \\
    $-u^2 P''$ & $-\sum_m a_m\,m(m-1)\,u^m$ \\
    $-2u\,P'$  & $-2\sum_m a_m\,m\,u^m$ \\
    $K\,P$     & $K\sum_m a_m\,u^m$ \\
    \bottomrule
  \end{tabular}
\end{table}

\begin{keybox}[El argumento clave]
Si la suma vale $0$ \textbf{para todo $u$}, entonces \textbf{cada coeficiente de
$u^m$ vale cero por separado}. Es el mismo razonamiento que con un polinomio: si
un polinomio es idénticamente nulo, todos sus coeficientes lo son.
\end{keybox}

Igualando a cero el coeficiente de $u^m$:

\[
a_{m+2}\,(m+2)(m+1) = a_m\big[m(m-1) + 2m - K\big] = a_m\big[m(m+1) - K\big]
\]

\[
\boxed{\;a_{m+2} = a_m\,\frac{m(m+1) - K}{(m+1)(m+2)}\;}
\]

Ésta es la \textbf{relación de recurrencia}, y \textbf{salta de dos en dos}.

\begin{figure}[H]
  \centering
  % Clase 7 — la relacion de recurrencia salta de 2 en 2: pares e impares forman
% cadenas independientes, y cuando K = n(n+1) la cadena se corta y queda un polinomio.
\begin{tikzpicture}[scale=1,
  coef/.style={draw=figblue, fill=figblue!8, rounded corners=2pt,
               minimum width=0.86cm, minimum height=0.56cm, font=\small},
  coefcero/.style={draw=figgray!60, fill=figgray!12, rounded corners=2pt,
               minimum width=0.86cm, minimum height=0.56cm, font=\small, text=figgray}]

% ---- cadena par ----
\node[coef] (a0) at (0,0.9)   {$a_0$};
\node[coef] (a2) at (2.0,0.9) {$a_2$};
\node[coef] (a4) at (4.0,0.9) {$a_4$};
\node[coefcero] (a6) at (6.0,0.9) {$a_6$};
\node[coefcero] (a8) at (8.0,0.9) {$a_8$};
\draw[figvec] (a0) -- (a2);
\draw[figvec] (a2) -- (a4);
\draw[figvec] (a4) -- (a6);
\draw[figvec] (a6) -- (a8);
\node[figlbl] at (-1.15,0.9) {pares};

% ---- cadena impar ----
\node[coefcero] (a1) at (0,-0.9)   {$a_1$};
\node[coefcero] (a3) at (2.0,-0.9) {$a_3$};
\node[coefcero] (a5) at (4.0,-0.9) {$a_5$};
\node[coefcero] (a7) at (6.0,-0.9) {$a_7$};
\node[coefcero] (a9) at (8.0,-0.9) {$a_9$};
\draw[figvecaux] (a1) -- (a3);
\draw[figvecaux] (a3) -- (a5);
\draw[figvecaux] (a5) -- (a7);
\draw[figvecaux] (a7) -- (a9);
\node[figlbl, text=figgray] at (-1.15,-0.9) {impares};

% el corte
\draw[figred, line width=1.1pt] (5.0,1.72) -- (5.0,0.08);
\node[figlbl, figred, align=center] at (5.0,2.25) {$K=n(n+1)$\\ corta ac\'a};

\node[figlblaux, align=center] at (4.0,-1.75)
  {Con $a_1=0$ toda la cadena impar es nula: se elige una paridad y se trabaja en ella.};

\node[figlbl, align=center] at (4.0,-2.7)
  {$a_{m+2}=a_m\,\dfrac{m(m+1)-K}{(m+1)(m+2)}$\ : el salto de dos deja dos cadenas que\\[0.3em]
   no se comunican, y el numerador se anula justo en $m=n$ cuando $K=n(n+1)$.};

\end{tikzpicture}

  \caption*{\small\itshape Las dos consecuencias del salto de dos, en una sola
  imagen: las cadenas no se comunican, y una de ellas se corta cuando
  $K=n(n+1)$.}
\end{figure}

\subsection{Primera observación: pares e impares no se hablan}

Como el salto es de 2, $a_0$ genera $a_2, a_4, a_6,\dots$ y $a_1$ genera
$a_3, a_5, a_7,\dots$, \textbf{sin mezclarse nunca}. El espacio de soluciones se
parte en dos subespacios que no se comunican: el de las funciones \textbf{pares}
y el de las \textbf{impares}.

\begin{notebox}[Es una clasificación, no una restricción]
Sin pérdida de generalidad se puede tomar $a_1=0$ (soluciones pares) o $a_0=0$
(soluciones impares). La solución general será una combinación lineal y por lo
tanto mezclará ambas, pero \textbf{como base se pueden separar}.
\end{notebox}

\subsection{Segunda observación: cuando la serie se corta}

Mirando el numerador de la recurrencia: si para algún entero $n$ ocurre que
$K = n(n+1)$, entonces al llegar a $m=n$ el numerador se anula, $a_{n+2}=0$, y
\textbf{todos los coeficientes siguientes son nulos}. La serie deja de ser
infinita: es un \textbf{polinomio}.

\[
\boxed{\;K = n(n+1)\quad\Longrightarrow\quad a_m = 0 \;\;\text{para todo } m > n\;}
\]

(unos se anulan por la recurrencia, otros por paridad).

\begin{keybox}[Y eso es una ganancia enorme]
En este caso no hay ninguna aproximación: la solución es \textbf{exacta}, y un
polinomio está definido en todo $\mathbb{R}$ ---en particular en todo el intervalo
$[-1,1]$ que interesa, extremos incluidos---. «Ahí ya no a lo físico o ingeniero,
ahí posta a lo matemático.»
\end{keybox}

\section{Los polinomios de Legendre}

\subsection{Cálculo de los primeros}

Todos salen de aplicar la recurrencia con $K = n(n+1)$. Nótese que \textbf{la
solución tiene siempre la misma paridad que $n$}.

\begin{table}[H]
  \centering
  \begin{tabularx}{\textwidth}{@{}c c X l@{}}
    \toprule
    $n$ & $K$ & \textbf{Recurrencia} & \textbf{Solución} \\
    \midrule
    0 & 0  & $a_2 = 0$ & $P_0(u) = a_0$ (constante) \\[0.3em]
    1 & 2  & $a_3 = 0$ & $P_1(u) = a_1\,u$ \\[0.3em]
    2 & 6  & $a_2 = a_0\dfrac{0-6}{1\cdot2} = -3a_0$,\ \ $a_4=0$
        & $P_2(u) = a_0\,(1 - 3u^2)$ \\[0.7em]
    3 & 12 & $a_3 = a_1\dfrac{1\cdot2-12}{2\cdot3} = -\tfrac{5}{3}a_1$,\ \ $a_5=0$
        & $P_3(u) = a_1\left(u - \tfrac{5}{3}u^3\right)$ \\
    \bottomrule
  \end{tabularx}
\end{table}

\begin{notebox}[La constante]
La solución constante ($n=0$) es solución, aunque «no es tan obvio: hay que hacer
una cuentita para darse cuenta».
\end{notebox}

En general, para cada entero $n$ hay una solución polinomial de \textbf{grado
$n$} y de \textbf{la misma paridad que $n$}. Son infinitas: los \textbf{polinomios
de Legendre}.

\subsection{Normalización}

Como la ecuación es lineal y homogénea, cualquier múltiplo de una solución es
solución: todas las de arriba quedaron a menos de una constante multiplicativa
($a_0$ o $a_1$). Para fijarla, la literatura adopta la convención

\[
\boxed{\;P_n(1) = 1\;}
\]

Con ella, los primeros son

\[
P_0(u) = 1
\qquad
P_1(u) = u
\qquad
P_2(u) = \tfrac{1}{2}\left(3u^2 - 1\right)
\qquad
P_3(u) = \tfrac{1}{2}\left(5u^3 - 3u\right)
\]

\begin{figure}[H]
  \centering
  % Clase 7 — los cuatro primeros polinomios de Legendre en el intervalo que importa,
% con la normalizacion P_n(1)=1 que fija la convencion de la literatura.
\begin{tikzpicture}
\begin{axis}[emfig, width=10.5cm, height=6.6cm,
  xlabel={$u=\cos\theta$}, ylabel={$P_n(u)$},
  domain=-1:1, samples=120,
  xmin=-1.15, xmax=1.15, ymin=-1.35, ymax=1.55,
  xtick={-1,-0.5,0,0.5,1}, ytick={-1,-0.5,0,0.5,1},
  axis lines=middle,
  % OJO: `axis lines=middle` REDEFINE `every axis x label`, asi que el estilo
  % del rotulo tiene que ir DESPUES o queda centrado sobre la linea del eje
  % (y el `=` se lee como `\neq`).
  xlabel style={anchor=north west, yshift=-3pt},
  legend style={at={(0.5,-0.16)}, anchor=north, legend columns=4, draw=figgray!40}]
  \addplot[figblue]              {1};                    \addlegendentry{$P_0$}
  \addplot[figred]               {x};                    \addlegendentry{$P_1$}
  \addplot[figamber]             {0.5*(3*x^2-1)};        \addlegendentry{$P_2$}
  \addplot[figgray, dashed]      {0.5*(5*x^3-3*x)};      \addlegendentry{$P_3$}
  % la normalizacion: todos pasan por (1,1)
  \addplot[only marks, mark=*, mark size=1.6pt, figblue, forget plot]
    coordinates {(1,1)};
\end{axis}
\node[figlbl, align=center] at (5.2,-2.35)
  {Todos valen $1$ en $u=1$ (normalizaci\'on) y tienen la misma paridad que $n$.\\[0.2em]
   $P_n$ es de grado $n$: por eso son linealmente independientes.};
\end{tikzpicture}

  \caption*{\small\itshape Los cuatro primeros. Lo que la figura hace visible es
  la normalización ---el punto común en $u=1$--- y que cada uno tiene un grado
  distinto, que es de donde saldrá la independencia lineal.}
\end{figure}

\begin{notebox}[La comprobación es directa]
En $P_2$, la forma $a_0(1-3u^2)$ evaluada en $u=1$ da $-2a_0$, así que hay que
elegir $a_0 = -\tfrac12$.
\end{notebox}

\begin{notebox}[Sobre el rol de todo esto]
«Alguien puede decir: esto no tiene nada de física, es sólo matemática.» Y en
parte es cierto. Lo que hay que llevarse no son estas soluciones particulares sino
\textbf{el método}: separación de variables, reducción a ecuaciones ordinarias, y
resolución de ésas por desarrollo en serie de potencias.
\end{notebox}

\section{Por qué las demás soluciones no sirven}

Falta el caso $K \neq n(n+1)$, en el que la serie \textbf{no se corta nunca}.

\subsection{Convergencia por el criterio del cociente}

Se aplica el \textbf{criterio de D'Alembert} (del cociente). La serie no es a
priori de términos positivos porque $u$ puede ser negativo, pero basta estudiar la
de los valores absolutos; y de hecho, como las soluciones son puras pares o puras
impares, sacando el primer término de factor sólo quedan potencias de $u^2$.

\[
\left|\frac{a_{m+2}\,u^{m+2}}{a_m\,u^{m}}\right|
= \left|\frac{m(m+1)-K}{(m+1)(m+2)}\right|\,u^{2}
\;\xrightarrow[m\to\infty]{}\; u^{2}
\]

Por lo tanto la serie \textbf{converge absolutamente si $|u| < 1$}, y en
$u = \pm1$ el criterio \textbf{no decide}.

\begin{warnbox}[Y $u=\pm1$ es justamente donde hace falta]
$u=1$ corresponde a $\theta = 0$ y $u=-1$ a $\theta = \pi$, los dos polos, que
forman parte del dominio del problema.
\end{warnbox}

\begin{notebox}[Ejercicio dejado en clase]
Explícitamente fuera del parcial ---«háganlo porque es divertido»---: probar que si
$K \neq n(n+1)$ con $n$ entero, la serie \textbf{no converge} en $u = \pm 1$. La
divergencia es \textbf{logarítmica}.
\end{notebox}

\begin{keybox}[Conclusión]
Sólo sirven las \textbf{soluciones polinomiales}. Es una divergencia afortunada:
todas las soluciones no polinomiales están mal definidas justo en el borde, así
que se pueden descartar. Lo que empezó como una restricción matemática ($K$ no
puede ser cualquier cosa) resulta ser una \textbf{condición física}.
\end{keybox}

\subsection{El caso \texorpdfstring{$K=0$}{K=0}, resuelto exactamente}

Para muestra vale un botón: se resuelve exactamente el caso más simple. Con $K=0$
y usando la forma compacta de la ecuación,

\[
\dv{u}\left[(1-u^2)\dv{P}{u}\right] = 0
\qquad\Longrightarrow\qquad
(1-u^2)\,P'(u) = A
\qquad\Longrightarrow\qquad
P'(u) = \frac{A}{1-u^2}
\]

Primitivando:

\[
\boxed{\;P(u) = \frac{A}{2}\ln\!\left(\frac{1+u}{1-u}\right) + B\;}
\]

\textbf{Verificación} de que la primitiva está bien:

\[
\dv{u}\left[\frac{A}{2}\ln\frac{1+u}{1-u}\right]
= \frac{A}{2}\left[\frac{1}{1+u} + \frac{1}{1-u}\right]
= \frac{A}{2}\cdot\frac{(1-u)+(1+u)}{1-u^2}
= \frac{A}{1-u^2}
\]

Esta solución general es una \textbf{combinación lineal de dos soluciones}: la
\textbf{constante} $B$, que es la polinomial ($P_0$), y el término logarítmico,
perfectamente definido en el intervalo abierto $(-1,1)$ pero que \textbf{explota
en $u=\pm1$}.

\begin{figure}[H]
  \centering
  % Clase 7 — el caso K=0 resuelto exacto: la solucion general es la constante
% (polinomial) mas un logaritmo que explota justo en u=+-1, que son los polos.
% Es el botón que muestra por que solo sirven las soluciones polinomiales.
\begin{tikzpicture}
\begin{axis}[emfig, width=10.5cm, height=6.6cm,
  xlabel={$u=\cos\theta$}, ylabel={$P(u)$},
  domain=-0.97:0.97, samples=200,
  xmin=-1.2, xmax=1.2, ymin=-2.8, ymax=2.8,
  xtick={-1,-0.5,0,0.5,1}, ytick={-2,-1,0,1,2},
  axis lines=middle,
  % OJO: `axis lines=middle` REDEFINE `every axis x label`, asi que el estilo
  % del rotulo tiene que ir DESPUES o queda centrado sobre la linea del eje
  % (y el `=` se lee como `\neq`).
  xlabel style={anchor=north west, yshift=-3pt},
  legend style={at={(0.5,-0.16)}, anchor=north, legend columns=2, draw=figgray!40}]
  \addplot[figblue]  {1};                          \addlegendentry{$P_0=1$ (polinomial)}
  \addplot[figred]   {0.5*ln((1+x)/(1-x))};        \addlegendentry{$\tfrac12\ln\frac{1+u}{1-u}$}
  % asintotas en los polos
  \addplot[figgray, dashed, forget plot] coordinates {(-1,-2.8) (-1,2.8)};
  \addplot[figgray, dashed, forget plot] coordinates {(1,-2.8) (1,2.8)};
\end{axis}
\node[figlblaux, text=figgray, fill=white, inner sep=1.5pt] at (0.62,5.35) {$\theta=\pi$};
\node[figlblaux, text=figgray, fill=white, inner sep=1.5pt] at (9.75,5.35) {$\theta=0$};
\node[figlbl, align=center] at (5.2,-2.35)
  {La divergencia cae exactamente en los dos polos, que \emph{s\'i} pertenecen al\\[0.2em]
   dominio del problema. Por eso las soluciones no polinomiales se descartan.};
\end{tikzpicture}

  \caption*{\small\itshape El caso exhibe en concreto lo que pasa en general: las
  soluciones no polinomiales están bien definidas en el interior y revientan
  exactamente donde el problema las necesita.}
\end{figure}

\section{La ecuación radial y la base de soluciones}

Con la angular resuelta, se vuelve a la radial, ya sabiendo que $K = n(n+1)$:

\[
\dv{r}\big(r^2 Z'(r)\big) = n(n+1)\,Z(r)
\]

Otra vez es lineal, de segundo orden, con coeficientes que dependen de la
variable, pero \textbf{ésta no es difícil}.

\begin{notebox}[Ejercicio con pista]
Hacer el cambio de variable $v = \ln r$.
\end{notebox}

El resultado es una combinación de dos potencias:

\[
\boxed{\;Z_n(r) = A_n\,r^{\,n} + B_n\,r^{-(n+1)}\;}
\]

Juntando ambas piezas, la \textbf{base de soluciones} de la ecuación de Laplace en
coordenadas esféricas con simetría azimutal es:

\begin{table}[H]
  \centering
  \begin{tabular}{@{}cl@{}}
    \toprule
    $n$ & \textbf{Soluciones (dos por cada $n$)} \\
    \midrule
    0 & $1$ \quad;\quad $\dfrac{1}{r}$ \\[0.6em]
    1 & $r\cos\theta$ \quad;\quad $\dfrac{\cos\theta}{r^{2}}$ \\[0.6em]
    2 & $r^{2}P_2(\cos\theta)$ \quad;\quad $\dfrac{P_2(\cos\theta)}{r^{3}}$ \\[0.6em]
    $n$ & $r^{\,n}P_n(\cos\theta)$ \quad;\quad $\dfrac{P_n(\cos\theta)}{r^{\,n+1}}$ \\
    \bottomrule
  \end{tabular}
\end{table}

\[
\boxed{\;\phi(r,\theta) = \sum_{n=0}^{\infty}\left[A_n\,r^{\,n} + \frac{C_n}{r^{\,n+1}}\right]P_n(\cos\theta)\;}
\]

\begin{keybox}[Lo que hay que llevarse del método]
Elegir las coordenadas que convengan al problema, proponer soluciones producto, y
así \textbf{reducir una ecuación en derivadas parciales a ecuaciones ordinarias}.
Que resolver esas ordinarias haya dado trabajo es otra cuestión; el problema
original era más difícil.
\end{keybox}

\section{Ejemplo: esfera conductora en un campo uniforme}

\subsection{Planteo y condiciones de borde}

Una \textbf{esfera conductora descargada} de radio $a$, colocada en un campo
eléctrico uniforme $\vec E_0 = E_0\hat z$. Se quiere el potencial para $r \ge a$.

\begin{figure}[H]
  \centering
  % Clase 7 — el problema resuelto: esfera conductora descargada de radio a en un
% campo uniforme. Se marcan las DOS condiciones de borde y el dato extra (Q=0)
% que hizo falta para cerrar la constante.
% Las dos CB van SIMETRICAS a media altura y con relleno blanco: apiladas arriba
% se comen el arco de theta y el rotulo de r (pasa en el render, no en el codigo).
\begin{tikzpicture}[scale=1]

  \draw[figaxis] (0,-2.5) -- (0,3.4);
  \node[figlblaux] at (0,3.62) {$z$};

  \foreach \xx in {-3.5,-2.85,2.85,3.5}{
    \draw[figvecaux] (\xx,-2.3) -- (\xx,2.9);
  }
  \node[figlblaux] at (-3.5,3.22) {$\vec E_0$};

  % esfera conductora
  \fill[figgray!25] (0,0) circle (1.0);
  \draw[figline, line width=1.1pt] (0,0) circle (1.0);
  \node[figlblaux] at (0,-0.36) {$\vec E=0$};
  \draw[figvecaux] (0,0) -- (218:1.0);
  \node[figlblaux, fill=figgray!25, inner sep=1.2pt] at (-0.4,-0.06) {$a$};
  \foreach \ang in {62,76,90,104,118}{\node[figpos] at (\ang:1.0) {};}
  \foreach \ang in {242,256,270,284,298}{\node[figneg] at (\ang:1.0) {};}

  % punto de observacion, bien arriba, con r y theta libres de las CB
  \coordinate (P) at (1.15,2.55);
  \node[figneutra] at (P) {};
  \draw[figvec] (0,0) -- (P);
  \node[figlbl, fill=white, inner sep=1.5pt] at (0.98,1.72) {$r$};
  \draw[figgray, line width=0.6pt] (0,1.95) arc (90:65.7:1.95);
  \node[figlblaux, fill=white, inner sep=1.2pt] at (0.3,2.12) {$\theta$};

  % las dos condiciones de borde, simetricas y a media altura
  \node[figlbl, figred, fill=white, inner sep=2pt] at (-2.55,0.9) {\textbf{CB 1}};
  \node[figlblaux, align=center, fill=white, inner sep=2pt] at (-2.55,0.28)
    {$\phi(a,\theta)=\phi_0$\\ (equipotencial)};
  \draw[figvecaux] (-1.78,0.5) -- (-0.88,0.42);

  \node[figlbl, figred, fill=white, inner sep=2pt] at (2.75,0.9) {\textbf{CB 2}};
  \node[figlblaux, align=center, fill=white, inner sep=2pt] at (2.75,0.28)
    {$\phi\to-E_0\,r\cos\theta$\\ cuando $r\to\infty$};

  % el dato extra
  \node[figlbl, align=center] at (0,-3.2)
    {\textbf{Dato extra que hace falta:} la esfera est\'a \emph{descargada}, $Q=0$.};
  \node[figlblaux, align=center] at (0,-3.95)
    {Sin \'el $\phi_0$ queda libre: cierran el problema las dos CB\\[0.15em]
     \emph{m\'as} el desarrollo multipolar (el t\'ermino en $1/r$ vale $Q/4\pi\varepsilon_0$).};

  \node[figlbl, align=center] at (0,-4.95)
    {$\phi(r,\theta)=-E_0\left(r-\dfrac{a^{3}}{r^{2}}\right)\cos\theta
     \qquad (r\ge a)$};

\end{tikzpicture}

  \caption*{\small\itshape Las dos condiciones de borde no alcanzan: el dato de
  que la esfera está descargada entra por una vía distinta, el desarrollo
  multipolar de las clases 4 y 5.}
\end{figure}

\begin{notebox}[Por qué vale Laplace]
Dentro de la esfera el campo es nulo (conductor en equilibrio) y el potencial es
constante: la solución interior es trivial y no interesa. Y en el exterior no hay
carga volumétrica.

El problema tiene simetría azimutal ---girar alrededor de $Oz$ no cambia nada---,
así que $\phi = \phi(r,\theta)$ y se aplica la base recién construida.
\end{notebox}

\textbf{Condición 1 --- en la superficie.} Un conductor en equilibrio es
\textbf{equipotencial}: todos sus puntos tienen el mismo potencial. Entonces

\[
\phi(a,\theta) = \phi_0 \qquad\text{para todo }\theta
\]

\textbf{Condición 2 --- en el infinito.} Lejos de la esfera su efecto desaparece y
se recupera el campo uniforme. Como el campo uniforme corresponde a un potencial
que varía \textbf{linealmente}, y $\vec E_0$ va según $\hat z$:

\[
\phi \;\xrightarrow[r\to\infty]{}\; -E_0\,z = -E_0\,r\cos\theta
\]

usando $z = r\cos\theta$.

\begin{notebox}[Sobre el signo menos]
Como $\vec E = -\grad\phi$, un campo que apunta según $+\hat z$ significa
potencial decreciente hacia $+z$. Físicamente: las cargas que crean el campo son
positivas de un lado y negativas del otro, y el potencial va de más a menos.
\end{notebox}

\subsection{Imponer la condición en el infinito}

En la serie general, cuando $r\to\infty$, todos los términos $C_n/r^{n+1}$
\textbf{tienden a cero} y de los términos $A_n r^n$ domina el de mayor $n$. Pero
el comportamiento pedido es \textbf{lineal en $r$}. Luego ningún término puede
crecer más rápido:

\[
A_n = 0 \qquad\text{para todo } n \ge 2
\]

Queda el término $n=1$, con $P_1(\cos\theta) = \cos\theta$:

\[
A_1\,r\cos\theta \;\equiv\; -E_0\,r\cos\theta
\qquad\Longrightarrow\qquad
\boxed{\;A_1 = -E_0\;}
\]

\begin{notebox}[Coincide exactamente en la forma]
Y no es casualidad: el potencial de un campo uniforme \emph{es} el elemento $n=1$
de la base.
\end{notebox}

Queda $A_0$, que multiplica a $P_0 = 1$: es una \textbf{constante aditiva}. No la
fija ninguna condición de borde porque el potencial está definido a menos de una
constante. \textbf{Se elige} $A_0 = 0$ ---una opción, no un dato---. El potencial
va quedando

\[
\phi(r,\theta) = -E_0\,r\cos\theta + \sum_{n=0}^{\infty}\frac{C_n}{r^{\,n+1}}P_n(\cos\theta)
\]

\subsection{Imponer la condición en la superficie}

Evaluando en $r=a$:

\[
-E_0\,a\cos\theta + \sum_{n=0}^{\infty}\frac{C_n}{a^{\,n+1}}P_n(\cos\theta) = \phi_0
\qquad\text{para todo }\theta
\]

\textbf{El argumento decisivo es de álgebra lineal:} los $P_n$ son polinomios
\textbf{todos de grado distinto}, y por lo tanto son \textbf{linealmente
independientes}. Así que se pueden identificar los coeficientes término a
término, escribiendo cada lado como desarrollo en la base $\{P_n\}$ ---no en la
base de monomios---.

\begin{warnbox}[Esta parte generó discusión en clase y vale precisarla]
El punto \textbf{no} es que los $P_n$ sean monomios (no lo son: $P_2$ tiene
término constante). El punto es que, teniendo grados distintos, forman una
\textbf{base del espacio de polinomios}, distinta de $\{1,u,u^2,\dots\}$ pero
igual de válida. Si una combinación lineal de ellos es idénticamente nula, se
anula primero el coeficiente del de mayor grado, luego el siguiente, y así hacia
abajo --- todos los coeficientes son cero.
\end{warnbox}

Identificando, con $\phi_0 = \phi_0\,P_0$ y $\cos\theta = P_1(\cos\theta)$:

\begin{table}[H]
  \centering
  \begin{tabular}{@{}lll@{}}
    \toprule
    \textbf{Orden} & \textbf{Ecuación} & \textbf{Resultado} \\
    \midrule
    $P_0$ & $\dfrac{C_0}{a} = \phi_0$ & $C_0 = \phi_0\,a$ \\[0.7em]
    $P_1$ & $-E_0\,a + \dfrac{C_1}{a^{2}} = 0$ & $C_1 = E_0\,a^{3}$ \\[0.7em]
    $P_n$, $n\ge2$ & $\dfrac{C_n}{a^{\,n+1}} = 0$ & $C_n = 0$ \\
    \bottomrule
  \end{tabular}
\end{table}

De una base infinita sobrevivieron \textbf{dos términos}. El docente lo llama, con
razón, una «solución de juguete».

\subsection{El dato que faltaba: la esfera está descargada}

Con las dos condiciones impuestas, \textbf{$\phi_0$ sigue sin determinarse}. Falta
usar un dato del enunciado que todavía no se usó: la esfera está
\textbf{descargada}.

El argumento es a través del \textbf{desarrollo multipolar} de las clases 4 y 5.
Lejos de la esfera hay dos contribuciones al campo: la del campo externo
$\vec E_0$, ya incorporada, y la de las cargas inducidas sobre la esfera. Esas
cargas forman una \textbf{distribución localizada}, así que su potencial admite
desarrollo multipolar, cuyo primer término va como $Q/4\pi\varepsilon_0 r$.

Pero la esfera está \textbf{polarizada, no cargada}: hay carga negativa de un lado
y positiva del otro, con \textbf{carga total nula}. Entonces el término en $1/r$
tiene que anularse, y ese término es exactamente $C_0/r$:

\[
C_0 = \phi_0\,a = 0 \qquad\Longrightarrow\qquad \phi_0 = 0
\]

\[
\boxed{\;\phi(r,\theta) = -E_0\left(r - \frac{a^{3}}{r^{2}}\right)\cos\theta
\qquad (r \ge a)\;}
\]

\begin{notebox}[Verificación inmediata]
En $r=a$ el paréntesis se anula y $\phi = 0$ para todo $\theta$, que es
exactamente la condición de equipotencial. Y para $r\to\infty$ el segundo término
desaparece y queda $-E_0 r\cos\theta$, el campo uniforme.

\textbf{Si la esfera estuviera cargada} con carga $Q$, el razonamiento sería el
mismo pero el coeficiente del término $1/r$ no sería nulo: valdría
$Q/4\pi\varepsilon_0$.
\end{notebox}

\begin{keybox}[Por qué esto importa]
Es el \textbf{primer problema del curso que no se podía resolver con las
herramientas de Física III}: la posición de las cargas era parte de la incógnita y
aun así el método la determinó. Es, en palabras del docente, «algo
cualitativamente más poderoso» --- aunque el caso concreto haya salido fácil.
\end{keybox}

\bigskip
\noindent\emph{Continúa en la Clase 8, calculando el campo eléctrico y la carga de
polarización inducida sobre la esfera, y pasando después a la misma técnica en
coordenadas cilíndricas.}

\end{document}
